\documentclass[aps,prx,nofootinbib,superscriptaddress,12pt]{revtex4-2}

\usepackage{graphicx}
\usepackage{hyperref}
\usepackage{amsmath, amssymb}
\usepackage{booktabs}

%% Bundle D4 — Topological quantum computation foundations.
%% Stage-10 FULL REDRAFT, 2026-08-15 (branch redraft/D4).
%%
%% Venue decision (this redraft): PRX Quantum, not Comm. Math. Phys.
%%   bundle_metadata.json records "Comm. Math. Phys. | PRX Quantum" and the
%%   charter says "or split between". A spine cannot be written against a pipe:
%%   CMP wants the categorical apparatus first, PRX Quantum wants the
%%   universality/compilation result first, and those orders are reverses.
%%   Measured inventory decides it — the only STRONG, novel, kernel-pure results
%%   in this bundle are the Fibonacci density theorem and its quantitative
%%   compilation strengthening. The categorical material is moderate identities
%%   plus two `rfl` encodings plus 18 native_decide apexes, which is not a CMP
%%   mathematics contribution. See papers/AutomatedReviews/2026-08-15-d4-stage10-redraft/D4.md.
%%
%% ORDERING NOTE — read before editing the section order.
%%   The section order is by ARGUMENT, not by subject-matter taxonomy, and it is
%%   deliberately NOT the "Onsager -> ... -> Fibonacci universality" chain the
%%   charter describes. Measured: the density theorem's project-side dependency
%%   set touches five modules (BraidGroup, FKLW.FibSU2Rep,
%%   FKLW.AharonovAradBridgeIteration, FKLW.OneParameterSubgroupSU2,
%%   FKLW.SU2BCHBracketClosure) and consumes NOTHING from Uqsl2, SU2kSMatrix,
%%   IsingBraiding, WRTInvariant, DrinfeldDouble or S3CenterAnyons. There is no
%%   chain; there is a set. Section 1 says so in plain words.

%% ═══════════════════════════════════════════════════════════════════
%% \figuredeferred{id}{reason} — the declared-deferral form
%% ═══════════════════════════════════════════════════════════════════
%%
%% WAVE_EXECUTION_PIPELINE.md §Stage 9: "A figure that is planned but not yet
%% drawn is declared with \figuredeferred{id}{reason}, never omitted silently."
%% The law mandated this form and, measured 2026-08-09, it had ZERO uses across
%% all 21 bundle drafts — nine of which carried no figures at all. Nine bundles
%% with no plan is a charter defect; nine bundles with a *stated* plan and a
%% stated blocker is an honest manuscript in progress, and a referee can tell
%% the two apart only if the plan is on the page.
%%
%% The deferral is DELIBERATELY VISIBLE in the compiled PDF. A silent macro
%% would satisfy `bundle_figure_adequacy` while showing a reader nothing, which
%% is the vacuous-Stage-9 failure ("ALL figures PASS" over an empty set) rebuilt
%% one layer up.
%%
%% ⚠️ The reason must name what BLOCKS the figure, not that it is absent.
%% "not yet drawn" is boilerplate and defeats the purpose; "awaiting the L=8
%% m→0 production run" is a blocker a reader can check and a later pass can
%% discharge.
%%
%% Definition lives here, beside docs/counts.tex, because 21 drafts consume it
%% and a per-draft copy is 21 places for the rendering to drift.

%% Rendered as a non-floating box on purpose. A `figure` float would be subject
%% to revtex's two-column placement rules across eleven drafts for no gain, and
%% a deferral does not need a figure number: it is a promise, not a result.

\providecommand{\figuredeferred}[2]{%
  \par\medskip\noindent
  \fbox{\parbox{0.92\columnwidth}{%
    \small\textbf{Figure deferred: \texttt{#1}.}
    Planned for this section; not yet produced. \emph{Blocked on:} #2.%
  }}%
  \par\medskip
}


\begin{document}

\title{Machine-checked universality of the Fibonacci braid representation,\\
with a graded verification of the surrounding modular-tensor-category apparatus}

\author{John Roehm}

\date{\today}

\begin{abstract}
We give a machine-checked proof, in Lean~4 and with no unverified
hypotheses, that the three-strand Fibonacci braid representation
$\rho_{\mathrm{Fib}}\colon B_3 \to \mathrm{SU}(2)$ has dense image: for
every $U \in \mathrm{SU}(2)$ and every $\varepsilon > 0$ there is a braid
word whose image agrees with $U$ entrywise to within $\varepsilon$. We
then strengthen density to efficiency, proving that the Dawson--Nielsen
compilation of a target $U$ returns, \emph{at one and the same
algorithmic level}, both an approximation of error at most $\varepsilon$
and a braid word of length at most
$C\,(\log 1/\varepsilon)^{c}$ with $c = \log 5 / \log(3/2)$ --- a
bundling that the existence-only statements of the Solovay--Kitaev
theorem do not provide. Around this result we develop the
modular-tensor-category apparatus it sits in: the Chevalley--Serre and
Hopf structure of $U_q(\mathfrak{sl}_2)$ and $U_q(\mathfrak{sl}_3)$;
$\mathrm{SU}(2)_k$ fusion with $S$-matrix unitarity and the Verlinde
formula verified entry-by-entry for $k \leq 5$; the Ising category's six
hexagon and four ribbon equations in $\mathbb{Q}(\zeta_{16})$, with the
closed-form invariants $\langle \mathrm{trefoil}\rangle = -\sqrt{2}$ and
$\mathrm{RT}(4_1) = -1/\sqrt{2}$; the Drinfeld centre
$\mathcal{Z}(\mathrm{Vec}_{S_3}) \simeq \mathrm{Rep}(D(S_3))$ with its
eight-object non-abelian spectrum; and the entanglement entropy of a
modular tensor category, where we prove that Casini--Huerta saturation
holds exactly when the category is trivial and that the deficit on a
non-trivial category is the topological entropy $\log\mathcal{D}$.
Because these layers were verified to different standards, we grade
them: each section states the strength of its own results, a summary
table appears in Sec.~\ref{sec:intro}, and a trust-base note records the
axiom set, the eighteen results discharged by compiler-level evaluation
rather than by the kernel, and the two hypotheses that remain assumed.
The layers are independent developments about a common subject rather
than a single derivation; we say where each stands and do not claim a
chain that the formalisation does not contain.
\end{abstract}

\maketitle

\tableofcontents

%% =====================================================================
%% §1 — Introduction
%% =====================================================================
\section{Introduction}
\label{sec:intro}

The Fibonacci anyon is the standard existence proof that topological
quantum computation is universal: braiding alone, with no measurement
and no additional gate, generates a dense subgroup of the unitary group
acting on the fusion space, so any target gate can be approximated to
arbitrary precision by a braid~\cite{FreedmanLarsenWang2002,
NayakSimonSternFreedmanDasSarma2008}. The argument is standard and it is
believed. What has not existed is a proof of it that a machine has
checked.

This paper supplies one. Our main theorem is that the three-strand
Fibonacci braid representation
$\rho_{\mathrm{Fib}} \colon B_3 \to \mathrm{SU}(2)$ has dense image, in
the concrete entrywise sense: for every $U$ in $\mathrm{SU}(2)$ and every
$\varepsilon > 0$ there exists a braid $b \in B_3$ with
$\lVert \rho_{\mathrm{Fib}}(b)_{ij} - U_{ij}\rVert < \varepsilon$ for all
$i,j$. The statement is unconditional --- it depends on no assumed
lemma, no tracked hypothesis and no numerical oracle --- and it is
checked by the Lean~4 kernel from the axioms of the ambient type theory.
Section~\ref{sec:density} proves it.

Density is an existence statement, and on its own it is not yet a
computational one: it does not say how long the braid is. The
Solovay--Kitaev theorem supplies the missing bound, and the
Dawson--Nielsen analysis makes it
quantitative~\cite{KitaevShenVyalyi2002,DawsonNielsen2006}. What we
verify in Sec.~\ref{sec:sk} is stronger than either in one specific
respect: our compilation theorem returns the error bound and the length
bound \emph{about the same compiled object at the same recursion depth},
rather than asserting that a short word exists and separately that an
accurate word exists. Existence-only forms leave open, as a matter of
logical form rather than of mathematical fact, whether the short word
and the accurate word are the same word. Ours does not.

\paragraph{What this paper is, and how to read its claims.}
The Fibonacci result sits inside a larger Lean development covering
quantum groups, fusion categories, modular data, knot invariants and
categorical entropy, and we present that development here. But those
layers were built at different times and to different standards, and
some of them are considerably weaker than their names suggest. Rather
than present a uniform assertion of verification and let a reader
discover the variation, we grade every layer explicitly, in the section
that carries it, against the scale in Table~\ref{tab:grades}.

\begin{table}[t]
\caption{Strength grading used throughout. Every section states the
grade of its own results. \emph{Kernel-checked} means the Lean~4 kernel
verified the proof term against the ambient axioms
$\{\texttt{propext}, \texttt{Classical.choice}, \texttt{Quot.sound}\}$;
\emph{evaluation-checked} means the proof invokes \texttt{native\_decide},
which delegates a decidable computation to compiled code and adds a
per-theorem axiom (Sec.~\ref{sec:ising}, Sec.~\ref{sec:trust}).}
\label{tab:grades}
\begin{ruledtabular}
\begin{tabular}{ll}
Grade & Meaning \\
\colrule
Theorem & A substantive statement, kernel-checked, no assumed hypothesis. \\
Theorem (evaluated) & Substantive and correct, but discharged by
  \texttt{native\_decide}; the trust base is the Lean compiler, not the kernel. \\
Conditional & Substantive, kernel-checked, but consumes a hypothesis
  the project has not discharged. \\
Instance & Verified for named finite cases, not for a general object. \\
Encoding & The statement unfolds to the definition it names; it fixes a
  convention and proves nothing about it. \\
\end{tabular}
\end{ruledtabular}
\end{table}

Two consequences of that grading are worth stating in the introduction
rather than burying.

First, the Witten--Reshetikhin--Turaev partition functions we record in
Sec.~\ref{sec:ising} are \emph{Encoding}-grade: the two headline
identities are true by unfolding the definitions that name them, and
they establish the normalisation convention rather than any invariance
property. Kirby-move invariance, which is the content a reader would
expect from a WRT section, is not proved here.

Second, the layer that relates horizon modular data to quantum error
correction --- which earlier project documents described as a
load-bearing biconditional --- is, measured against its Lean statements,
a set of definitions ($\textsf{codeDistance} := \log d_{\max}$,
$\textsf{scramblingTimeBound} := \log \mathcal{D}^2$) whose theorems are
the elementary theory of the logarithm on $[1,\infty)$. It carries no
error-correcting code, no recovery map and no scrambling dynamics.
Section~\ref{sec:entropy} says so in those words and confines it to a
scope paragraph. The genuinely categorical entropy results --- the
Casini--Huerta saturation biconditional and the $\log\mathcal{D}$
deficit --- are separate theorems, they do contain the category, and
they are what that section is built on.

\paragraph{The layers are independent, not a chain.}
It would be natural to present this development as a single verified
chain running from a quantum group up to universality, and earlier
framings of this project did. That description is not accurate and we do
not use it. The density theorem of Sec.~\ref{sec:density} is proved from
an explicit pair of $2\times 2$ matrices satisfying the Yang--Baxter
relation together with a Cartan-subgroup argument; its project-side
dependencies touch the braid-group and $\mathrm{SU}(2)$ analysis modules
and nothing else. It does not consume the $U_q(\mathfrak{sl}_2)$
presentation of Sec.~\ref{sec:qgroups}, the modular data of
Sec.~\ref{sec:modular}, the Ising category of Sec.~\ref{sec:ising}, or
the Drinfeld centre of Sec.~\ref{sec:drinfeld}. Those sections are
independent formalisations about a common subject. Presenting them as
prerequisites for the main theorem would misrepresent what was proved,
and a reader inspecting the sources would find it out.

\paragraph{Relation to the companion compilation paper.}
The alphabet-agnostic and dimension-agnostic theory of quantitative
gate compilation --- the generic \textsf{GeneratingSet} abstraction, the
lift to $\mathrm{SU}(d)$, and the instantiations at Clifford$+T$,
Read--Rezayi $\mathrm{SU}(2)_5$ and $\mathrm{SU}(2)_7$, trapped-ion
$\mathrm{SU}(4)$ and Clifford$+\mathrm{CCZ}$ $\mathrm{SU}(8)$ --- is
developed in a companion paper~\cite{RoehmD8Compilation} and is not
reproduced here. This paper owns the Fibonacci alphabet, which is the
anyonic one; the companion paper owns the general theory and cites this
one for the Fibonacci anchor. Section~\ref{sec:sk} states the boundary
precisely.

\paragraph{Roadmap.}
Section~\ref{sec:density} proves density.
Section~\ref{sec:sk} strengthens it to a quantitative compilation bound.
Section~\ref{sec:gate} gives an independent gate-level result: a
two-qubit dark state with a sign flip on the $\pi$ loop and vanishing
dynamical phase.
Sections~\ref{sec:qgroups}--\ref{sec:drinfeld} develop the algebraic
apparatus --- quantum groups, modular data and the Verlinde formula, the
Ising category and its knot invariants, and the Drinfeld centre.
Section~\ref{sec:entropy} treats the entanglement entropy of a modular
tensor category.
Section~\ref{sec:open} collects what remains open. A trust-base note
before the bibliography records the axiom set and the formalisation's
exact standing.

%% =====================================================================
%% §2 — Fibonacci density  [TO DRAFT]
%% =====================================================================
\section{The Fibonacci braid representation and its density in $\mathrm{SU}(2)$}
\label{sec:density}

%% TODO(redraft): draft this section. Content:
%%  - the representation: sigma_1 = diag(R_1, R_tau), sigma_2 = F sigma_1 F,
%%    det-normalised to SU(2); rho_Fib_SU2 = braidGroup3HomFromPair on the
%%    Yang-Baxter pair. STATE that the F and R data are POSITED, not derived
%%    from the SU(2)_3 fusion category in this development.
%%  - DenseInSpecialUnitary 3 2 rho unfolds to the entrywise density statement.
%%  - the proof shape: Cartan closed-subgroup argument; discrete branch excluded
%%    by an Aharonov-Arad cardinality argument.
%%  - GRADE: Theorem (kernel-checked, unconditional).

%% =====================================================================
%% §3 — Quantitative Solovay-Kitaev  [TO DRAFT]
%% =====================================================================
\section{Quantitative Solovay--Kitaev at the Fibonacci alphabet}
\label{sec:sk}

%% TODO(redraft): draft this section. Four jobs only:
%%  1. the bundled statement and why existence-only is strictly weaker
%%  2. the exponent c = log5/log(3/2), with the two sentinels 3 < c < 4
%%  3. the calibration K_proof <= 788 against K_compose = 1024
%%  4. the D8 boundary, one paragraph
%% DO NOT refill with Path-A recursion internals or per-module LoC tables.

%% =====================================================================
%% §4 — SPT two-qubit gate  [TO DRAFT]
%% =====================================================================
\section{A symmetry-protected two-qubit gate substrate}
\label{sec:gate}

%% TODO(redraft): dark state in singlet-sector kernel, non-trivial;
%% sign flip on the pi loop AND vanishing dynamical phase (both PROVED —
%% the prior draft's "explicitly deferred" was stale). Holonomy-level
%% adiabatic transport is what genuinely remains open.

%% =====================================================================
%% §5 — Quantum groups  [TO DRAFT]
%% =====================================================================
\section{Quantum groups: the algebraic input}
\label{sec:qgroups}

%% TODO(redraft)

%% =====================================================================
%% §6 — Modular data  [TO DRAFT]
%% =====================================================================
\section{Modular data: fusion, $S$-matrix unitarity, and the Verlinde formula}
\label{sec:modular}

%% TODO(redraft)

%% =====================================================================
%% §7 — Ising category  [TO DRAFT]
%% =====================================================================
\section{The Ising category: braiding, knot invariants, and the trust boundary}
\label{sec:ising}

%% TODO(redraft). D^2 = 4, D = 2. NOT D = sqrt 2 (that is d_sigma).
%% WRT is a SUBSECTION here, three sentences, Encoding-grade.

%% =====================================================================
%% §8 — Drinfeld centre  [TO DRAFT]
%% =====================================================================
\section{The Drinfeld centre and a non-abelian anyon spectrum}
\label{sec:drinfeld}

%% TODO(redraft)

%% =====================================================================
%% §9 — Categorical entanglement entropy  [TO DRAFT]
%% =====================================================================
\section{Entanglement entropy of a modular tensor category}
\label{sec:entropy}

%% TODO(redraft)

%% =====================================================================
%% §10 — Open problems  [TO DRAFT]
%% =====================================================================
\section{Open problems}
\label{sec:open}

%% TODO(redraft)

%% =====================================================================
%% Trust base (unnumbered)  [TO DRAFT]
%% =====================================================================
\section*{Formalisation and trust base}
\label{sec:trust}

%% TODO(redraft)

\section*{Methods and tools disclosure}
The formal verification layer of this work was developed with
substantial machine assistance: Lean~4 proof drafting and iteration
used LLM-based agents (Anthropic Claude-family models) operating under
the project's staged validation pipeline. Every proof is ultimately
certified by the Lean~4 kernel via \texttt{lake build} on the public
repository, independently of any AI tool, with the explicit exception of
the \texttt{native\_decide} results itemised in
Sec.~\ref{sec:trust}. Manuscript prose was AI-assisted and
human-reviewed.

\begin{thebibliography}{99}

\bibitem{AharonovArad2011}
D.~Aharonov and I.~Arad,
``The BQP-hardness of approximating the Jones polynomial,''
New J.\ Phys.\ \textbf{13}, 035019 (2011),
arXiv:quant-ph/0605181.

\bibitem{CasiniHuerta2009}
H.~Casini and M.~Huerta,
``Entanglement entropy in free quantum field theory,''
J.\ Phys.\ A \textbf{42}, 504007 (2009),
arXiv:0905.2562.

\bibitem{DawsonNielsen2006}
C.~M.~Dawson and M.~A.~Nielsen,
``The Solovay--Kitaev algorithm,''
Quantum Inf.\ Comput.\ \textbf{6}, 81--95 (2006),
arXiv:quant-ph/0505030.

\bibitem{FreedmanLarsenWang2002}
M.~H.~Freedman, M.~J.~Larsen, and Z.~Wang,
``A modular functor which is universal for quantum computation,''
Commun.\ Math.\ Phys.\ \textbf{227}, 605 (2002),
arXiv:quant-ph/0001108.

\bibitem{KitaevAnyons2003}
A.~Yu.~Kitaev,
``Fault-tolerant quantum computation by anyons,''
Ann.\ Phys.\ \textbf{303}, 2 (2003),
arXiv:quant-ph/9707021.

\bibitem{KitaevPreskill2006}
A.~Kitaev and J.~Preskill,
``Topological entanglement entropy,''
Phys.\ Rev.\ Lett.\ \textbf{96}, 110404 (2006),
arXiv:hep-th/0510092.

\bibitem{KitaevShenVyalyi2002}
A.~Yu.~Kitaev, A.~H.~Shen, and M.~N.~Vyalyi,
\emph{Classical and Quantum Computation},
Graduate Studies in Mathematics \textbf{47}, American Mathematical
Society (2002).

\bibitem{Lickorish1997}
W.~B.~R.~Lickorish,
\emph{An Introduction to Knot Theory},
Graduate Texts in Mathematics \textbf{175}
(Springer, New York, 1997).

\bibitem{NayakSimonSternFreedmanDasSarma2008}
C.~Nayak, S.~H.~Simon, A.~Stern, M.~Freedman, and S.~Das~Sarma,
``Non-Abelian anyons and topological quantum computation,''
Rev.\ Mod.\ Phys.\ \textbf{80}, 1083 (2008),
arXiv:0707.1889.

\bibitem{ReshetikhinTuraev1991}
N.~Reshetikhin and V.~G.~Turaev,
``Invariants of $3$-manifolds via link polynomials and quantum groups,''
Invent.\ Math.\ \textbf{103}, 547 (1991).

\bibitem{RoehmD8Compilation}
J.~G.~Roehm,
``Kernel-verified universal quantum gate compilation:
alphabet-agnostic Solovay--Kitaev across dimensions,''
companion manuscript, in preparation (2026).

\bibitem{Verlinde1988}
E.~Verlinde,
``Fusion rules and modular transformations in $2$D conformal field
theory,''
Nucl.\ Phys.\ B \textbf{300}, 360 (1988).

\bibitem{Witten1989}
E.~Witten,
``Quantum field theory and the Jones polynomial,''
Commun.\ Math.\ Phys.\ \textbf{121}, 351 (1989).

\end{thebibliography}

\end{document}
